%%% document class and general options
\documentclass[%
pdftex,
11pt,
oneside,
headings=small,
ngerman,
parskip=half, % set to false for no indention
footsepline
]{scrartcl}

%%% page geometry
% variables
\newcommand{\papersize}[1]{\geometry{#1}}
\newcommand{\orientation}[1]{\geometry{#1}}

% geometry
\usepackage[
includehead,
includefoot,
left=20mm,
right=20mm,
top=10mm,
bottom=10mm,
%headsep=3.5\baselineskip,
headheight=55.83582pt,
]{geometry}

%%% page style
\usepackage{scrlayer-scrpage}

% style
\pagestyle{scrheadings}
\setkomafont{pageheadfoot}{\sffamily\bfseries\large}
\setkomafont{pagefoot}{\sffamily\mdseries\scriptsize}

% variables
\newcommand{\mytitle}{}
\newcommand{\mysubject}{}
\newcommand{\teacher}{}
\newcommand{\course}{}
\newcommand{\logo}{}
\newcommand{\icon}{}

% table handling
\usepackage{booktabs}
\usepackage{tabularx}
\usepackage{longtable}
\usepackage{ltablex}
\usepackage{tabu}
\usepackage{colortbl}
\usepackage{multirow}

% array stretch
\renewcommand{\arraystretch}{1.0}

% define new column types for the tabularx environment
\renewcommand{\tabularxcolumn}[1]{m{#1}}
\newcolumntype{Y}{>{\centering\arraybackslash}X}
\newcolumntype{Z}{>{\flushleft\arraybackslash}X}

% define new column types for the tabu environment
\newcolumntype{P}[1]{>{\centering\arraybackslash}p{#1}}
\newcolumntype{M}[1]{>{\centering\arraybackslash}m{#1}}
\newcolumntype{B}[1]{>{\centering\arraybackslash}b{#1}}

% head
\chead{%
%\tabulinesep=0pt%
%\extrarowsep=0pt%
%\begin{tabu} to \linewidth {X[1,c,m] M{153pt} @{\hspace{2.5pt}}}%
%  & \includegraphics[width=153pt]{\logo} \\
%\end{tabu}
\tabulinesep=2pt%
\extrarowsep=0pt%
\arrayrulewidth=0.5pt%
\begin{tabu} to \linewidth {|@{\hspace{1.5pt}}M{20pt}@{\hspace{1.5pt}}|X[1,c,m]|}
  \firsthline
  \includegraphics[width=20pt]{\icon} & \mytitle \\
  \lasthline
\end{tabu}}

% foot
\usepackage{totcount}
\usepackage{ifthen}
\usepackage{etoolbox}

\regtotcounter{page}

\cfoot{\ifnumcomp{\totvalue{page}}{>}{1}{Seite~\thepage}{}}
\ifblank{\teacher}
{%
  \ifoot{\mysubject}
  \ofoot{\course}
}{%
  \ifoot{\mysubject--\course}
  \ofoot{\teacher}
}%

%%% include additional packages
\usepackage[singlespacing]{setspace}
\usepackage{multicol}

% language and encodings
\usepackage[utf8]{inputenc}
\usepackage[ngerman]{babel}
\usepackage[T1]{fontenc}

% font
%\usepackage{droid}
\usepackage{helvet}
%\usepackage[default,osfigures,scale=0.95]{opensans}
\usepackage{lmodern}

% Times
%\usepackage{mathptmx}
%\usepackage[scaled=.92]{helvet}
%\usepackage{courier}

% Palatino
%\usepackage{mathpazo}
%\usepackage[scaled=.95]{helvet}
%\usepackage{courier}

% Utopia
%\usepackage{fourier}

\usepackage{eulervm}
\usepackage{upgreek}
\usepackage{eurosym}

% microtype
\usepackage[protrusion=true,expansion]{microtype}

% amsmath
\usepackage[sumlimits]{amsmath}
\usepackage{amsthm}
\usepackage{amssymb}
\usepackage{amsfonts}
\usepackage{mathabx}
\usepackage{relsize}
\usepackage{esvect}

% graphics
\usepackage[pdftex]{graphicx}
\graphicspath{{images/}}
\usepackage{subfig}

% portable graphics format
\usepackage{tikz}
\usetikzlibrary{arrows,decorations,backgrounds,fit,positioning,shapes,chains}
\usepackage{pgf}
\usepackage{pgfplots}
%\pgfplotsset{width=14cm, compat=1.10}
\pgfplotsset{compat=1.16}
\usepgfplotslibrary{fillbetween}

% control the placement of floats
\usepackage{flafter}
\usepackage[section]{placeins}


% handling of units
%\usepackage{units}
%\usepackage{nicefrac}
\usepackage{siunitx}
\sisetup{ output-decimal-marker = {,} }

% enumerate
\usepackage{enumerate}
\usepackage{enumitem}

%%% document wide settings
% enumeration label
\renewcommand{\labelenumi}{\textbf{\alph{enumi})}}
\setlist[enumerate]{leftmargin=*,itemsep=4pt,parsep=4pt,topsep=4pt,partopsep=4pt}

%%% command and environment definitions
%% box definitions
\usepackage{fancybox}

% circle arround text
\newcommand*\circled[1]{\tikz[baseline=(char.base)]{
  \node[shape=circle,draw,inner sep=2pt] (char) {#1};}}

% a box indicating another page
\newcommand{\npbox}[1][]{
  \vfill
  \hfill
  {\footnotesize \shadowbox{\colorbox{white}{Bitte wenden ...}}  \newpage }
}

% a box with shadow
\newcommand{\sframebox}[2][0.9\linewidth]{
  \begin{center}
    \vspace{3pt}
    \shadowbox{\colorbox{white}
      {\parbox[t]{#1-2\fboxsep-2\fboxrule-\shadowsize}{#2}}}
    \vspace{3pt}
  \end{center}
}

% a simple box
\newcommand{\simplebox}[2][0.9\linewidth]{
  \begin{center}
    \fbox{\parbox[t]{#1-2\fboxsep-2\fboxrule}{ #2 }}
  \end{center}
}

%
\newcommand{\bbox}[1]{
  \definecolor{structureshaded}{gray}{0.75}
  \definecolor{othercolor}{rgb}{1.0,1.0,1.0}
  \fcolorbox{structureshaded}{othercolor}{ #1 }
}

%
\newcommand{\cbox}[2][othercolor]{
  \definecolor{structureshaded}{gray}{0.75}
  \definecolor{othercolor}{rgb}{0.9,0.9,0.9}
  \begin{center}
    \fcolorbox{structureshaded}{#1}{\parbox[t]{1\linewidth}{#2}}
  \end{center}
}

%% list environment "aufgaben"
\newcounter{nummer}
\newenvironment{aufgaben}[1][1]
{\begin{list}{\ifthenelse{\value{nummer}=1}{\setcounter{nummer}{#1}}{}
  \normalsize{\textbf{Aufgabe~\arabic{nummer}~}}}{%
  \usecounter{nummer}
  \settowidth{\labelwidth}{\textbf{Aufgabe~\arabic{nummer}~~}}
  \setlength{\itemindent}{\labelwidth}
  \setlength{\leftmargin}{0em}
  \setlength{\leftmargini}{2em}
  \setlength{\parsep}{0.5ex plus0.2ex minus0.2ex}
  \setlength{\itemsep}{2ex}}}%
{\end{list}}%

%% tikz
%\begin{graph}[step size]{xmin}{xmax}{ymin}{ymax}
\newenvironment{graph}[5][1]{%
\begin{tikzpicture}[x=#1cm,y=#1cm]
  % grid
  \draw[very thin, color=lightgray, xstep=0.5, ystep=0.5]
    (#2-.4,#4-.4) grid (#3+.4,#5+.4);
  % axes
  \draw[->] (#2-.3,0) -- (#3+.3,0) node[above] {$x$};
  \draw[->] (0,#4-.3) -- (0,#5+.3) node[right] {$y$};
  % units for cartesian reference frame
  \foreach \x/\xtext in {#2,...,#3}
    \draw (\x cm,1pt) -- (\x cm,-3pt) node[anchor=north] {$\xtext$};
  \foreach \y/\ytext in {#4,...,#5}
    \draw (1pt,\y cm) -- (-3pt,\y cm) node[anchor=east] {$\ytext$};
}{\end{tikzpicture}}%

%\begin{tsdiagram}[step size]{xmin}{xmax}{ymin}{ymax}
\newenvironment{tsdiagram}[5][1]{%
\begin{tikzpicture}[x=#1cm,y=#1cm]
  % grid
  \draw[very thin, color=gray, xstep=0.5, ystep=0.5]
    (#2-.4,#4-.4) grid (#3+.4,#5+.4);
  % axes
  \draw[->] (#2-.3,0) -- (#3+.3,0) node[above] {$t$};
  \draw[->] (0,#4-.3) -- (0,#5+.3) node[right] {$s$};
  % units for cartesian reference frame
  \foreach \x/\xtext in {#2,...,#3}
    \draw (\x cm,1pt) -- (\x cm,-3pt) node[anchor=north] {};
  \foreach \y/\ytext in {#4,...,#5}
    \draw (1pt,\y cm) -- (-3pt,\y cm) node[anchor=east] {};
}{\end{tikzpicture}}%

%\function{xmin}{xmax}{gnuplot-input}{caption}
\newcounter{functionid}
\newcommand{\function}[4]
{%
  \draw[smooth, line width=0.6pt, samples=400, domain=#1:#2]
    plot[id=\arabic{functionid}] function{#3} node [right] {#4};
  \stepcounter{functionid}
}%

%%% math command definitions
\theoremstyle{plain}
\newtheorem{defn}{Definition}[section]
\newtheorem{lem}[defn]{Lemma}
\newtheorem{thm}[defn]{Satz}
\newtheorem{kor}[defn]{Korollar}
\newtheorem{bem}[defn]{Bemerkung}

\newcommand{\mbf}[1]{\ensuremath{\mathbf{#1}}}
\newcommand{\bs}[1]{\ensuremath{\boldsymbol{#1}}}

\renewcommand{\div}[1]{\ensuremath{\nabla\cdot\mbf{#1}}}
\newcommand{\rot}[1]{\ensuremath{\nabla\times\mbf{#1}}}
\newcommand{\rotrot}[1]{\ensuremath{\nabla\times\rot{#1}}}
\newcommand{\grad}[1]{\ensuremath{\nabla #1}}

\renewcommand{\arg}{\ensuremath{\left( \mbf{x} \right)}}
\newcommand{\args}{\ensuremath{\left( \mbf{x},t \right)}}

\newcommand{\dt}[1]{\ensuremath{\dfrac{\partial #1}{\partial t}}}
\newcommand{\ddt}[1]{\ensuremath{\dfrac{\partial^2 #1}{\partial t^2}}}
\newcommand{\dtb}[1]{\ensuremath{\dfrac{\partial\mbf{#1}}{\partial t}}}
\newcommand{\ddtb}[1]{\ensuremath{\dfrac{\partial^2\mbf{#1}}{\partial t^2}}}

\renewcommand{\sinh}[1]{\ensuremath{\mathop{\mathrm{sinh}}\!\left( #1 \right)}}
\renewcommand{\cosh}[1]{\ensuremath{\mathop{\mathrm{cosh}}\!\left( #1 \right)}}
\renewcommand{\tanh}[1]{\ensuremath{\mathop{\mathrm{tanh}}\!\left( #1 \right)}}
\renewcommand{\exp}[1]{\ensuremath{\mathop{\mathrm{e}^{#1}}}}

\newcommand{\dd}{\ensuremath{\mathrm{d}}}

\newcommand{\myunit}[1]{\;\unit{#1}}
\newcommand{\myunitfrac}[2]{\;\unitfrac{#1}{#2}}
\newcommand{\mybunit}[1]{\;\left( \unit{#1} \right)}
\newcommand{\mybunitfrac}[2]{\;\left( \unitfrac{#1}{#2} \right)}

\newcommand*{\molec}[3]{\textsf{#1\textsubscript{#2}\textsuperscript{#3}}}
